%% kapa: exam=1 round=20 year=2009 period=1 form=A subject=civil_law subject_ko=민법 seq=39 number=39
%% source: Q-Net 감정평가사 기출문제 (cst00309) artlSeq=5024178 file=extracted/1008/5024178/2009년 제20회 감정평가사 1차시험 기출문제/1교시.hwp
%% conversion: kordoc 4.24.51 -> tools/gen1.py (auto); review=proofread-sonnet; flags=none
\begin{question}{39}
17세 甲은 자신 소유의 부동산을 법정대리인의 동의 없이 乙에게 매도ㆍ등기하고, 乙은 다시 丙에게 매도ㆍ등기하였다. 다음 중 옳지 \uline{않은} 것은?
\begin{choices}
  \item 甲이 乙과의 매매계약을 취소한 경우에 丙은 선의ㆍ무과실이더라도 소유권을 취득할 수 없다.
  \item 甲은 乙과의 매매계약을 취소하고 이를 이유로 丙명의의 소유권이전등기의 말소를 청구할 수 있다.
  \item 甲이 매매계약을 취소하려면 丙이 아니라 乙에게 취소의 의사표시를 해야 한다.
  \item 甲이 乙과의 매매계약을 취소하여 甲이 부동산의 소유권을 회복한 경우에는 丙은 乙에게 매도인의 담보책임을 물을 수 있다.
  \item 만약 乙도 미성년자이며 법정대리인의 동의없이 甲 및 丙과의 매매계약을 체결한 경우, 甲의 법정대리인이 甲ㆍ乙간의 계약을 추인한 때에는 乙은 甲과의 계약을 취소할 수 없다.
\end{choices}
\end{question}
\answer{5}
